\section{Conclusion}
\label{sec:conclusion}
Phase-scoped QB offers a concrete way to constrain attention query geometry while retaining surrounding target batching. The original CUDA studies establish bitwise PQ/QB attention correspondence on a fixed replay and lower complete-request cost on the registered three-family panel. Later mapping experiments show that execution mode, reference admissibility and representation preparation constrain how far that local result can be optimized.

The newer B570 and RTX request panels add an important distinction. Both qualify real q=2 execution and preserve the registered endpoints and state traces, yet neither establishes a practical QB-over-PQ layout gain under its predeclared screen. Their speculative economics differ: the RTX policy is faster than target-only on its finite exposed panel, whereas the B570 policy is substantially more expensive. These are bounded results under different runtime contracts, not a cross-device ranking.

The evidence justifies an empirical architecture discussion about reference-preserving graph transformations, backend support and complete cost accounting. It does not yet justify a new accelerator design. A hardware direction should follow only if a measurable request-critical residual survives the strongest qualified software baseline and can be tied to a concrete physical mechanism.
